\solchapter{qm}

\sol{qm}{1} Electron: $\lambda = h/\sqrt{2mE} = 6.63\times10^{-34}/\sqrt{2(9.11\times10^{-31})(1.6\times10^{-17})} = 0.12\un{nm}$, comparable to the spacing of atoms in a crystal, so crystals diffract it. Baseball: $\lambda = h/mv = 1.1\times10^{-34}\un{m}$; a doorway diffracts it by an angle of $10^{-34}$ radians.

\sol{qm}{2} $\int A^2e^{-2\abs x/a}\dd x = A^2a = 1$, so $A = 1/\sqrt a$. $P(\abs x<a) = 1 - e^{-2} = 0.86$. $\langle x\rangle = 0$ by symmetry, and $\langle x^2\rangle = (2/a)\int_0^\infty x^2e^{-2x/a}\dd x = a^2/2$.

\sol{qm}{3} $E_n = n^2h^2/8mL^2$ with $L = 1\un{nm}$: $0.38$, $1.50$, and $3.38\un{eV}$. The $2\to1$ photon has $1.13\un{eV}$, wavelength $1240/1.13 = 1100\un{nm}$: infrared, invisible.

\sol{qm}{4} A $400\un{nm}$ photon carries $3.10\un{eV}$, so $K_{\max} = 0.8\un{eV}$. A $600\un{nm}$ photon carries $2.07\un{eV}<2.3\un{eV}$: no electrons at all, however bright the light.

\sol{qm}{5} $\dd E/\dd r = -\hbar^2/mr^3 + e^2/4\pi\epsilon_0r^2 = 0$ gives $r = 4\pi\epsilon_0\hbar^2/me^2 = a_0 = 0.053\un{nm}$ and $E = -e^2/8\pi\epsilon_0a_0 = -13.6\un{eV}$. With a muon, $m$ grows by $207$ (more precisely the reduced mass by $186$), so muonic hydrogen is about $200$ times smaller, $2.8\times10^{-13}\un{m}$, and bound by about $2.5\un{keV}$. Its size is close enough to the proton's that its energy levels are sensitive to the proton's radius --- which is how the proton radius was measured precisely in 2010.

\sol{qm}{6} $\kappa = \sqrt{2m(4\un{eV})}/\hbar = 1.0\times10^{10}\un{m^{-1}}$, so $T\approx e^{-2\kappa a} = e^{-10.2}\approx4\times10^{-5}$. Doubling the width squares this, to about $10^{-9}$. A proton has $\kappa$ larger by $\sqrt{1836} = 43$, giving $e^{-440}$: never.

\sol{qm}{7} $[\hat H,\hat a^\dagger] = \hbar\omega[\hat a^\dagger\hat a,\hat a^\dagger] = \hbar\omega\hat a^\dagger[\hat a,\hat a^\dagger] = \hbar\omega\hat a^\dagger$. So $\hat H\hat a^\dagger\ket n = \hat a^\dagger\hat H\ket n + \hbar\omega\hat a^\dagger\ket n = (E+\hbar\omega)\hat a^\dagger\ket n$. In position space, $\hat a\ket0 = 0$ reads $(x + \tfrac{\hbar}{m\omega}\tfrac{\dd}{\dd x})\psi_0 = 0$, whose solution is the Gaussian $\psi_0\propto e^{-m\omega x^2/2\hbar}$.

\sol{qm}{8} Each Pauli matrix has eigenvalues $\pm1$, with eigenvectors $\ket\uparrow,\ket\downarrow$ for $\sigma_z$; $(\ket\uparrow\pm\ket\downarrow)/\sqrt2$ for $\sigma_x$; and $(\ket\uparrow\pm i\ket\downarrow)/\sqrt2$ for $\sigma_y$. A spin up along $x$ gives $z$-results up or down with probability $\tfrac12$ each. Along an axis at angle $\theta$ to $x$, the probability of ``up'' is $\cos^2(\theta/2) = (1+\cos\theta)/2$. In the sequence $z\to x\to z$, the $x$ result is random, and the final $z$ result is again random, $50$--$50$, regardless of the first $z$ result: the $x$ measurement erased it.

\sol{qm}{9} $[\hat x,\hat p]f = x(-i\hbar f') + i\hbar(xf)' = i\hbar f$. For angular momentum, the only non-commuting pieces of $[yp_z - zp_y,\,zp_x - xp_z]$ are $[yp_z, zp_x] = yp_x[p_z,z] = -i\hbar yp_x$ and $[zp_y, xp_z] = xp_y[z,p_z] = i\hbar xp_y$. Together, $[\hat L_x,\hat L_y] = i\hbar(xp_y - yp_x) = i\hbar\hat L_z$.

\sol{qm}{10} $\Delta E = \lambda\langle x^4\rangle = 3\lambda\hbar^2/4m^2\omega^2$. First-order perturbation theory is trustworthy when this is much smaller than the level spacing $\hbar\omega$, i.e.\ when the dimensionless coupling $\lambda\hbar/m^2\omega^3\ll1$. (Even then, the full perturbation series in $\lambda$ diverges --- it is only asymptotic --- because for $\lambda<0$ the potential has no bottom.)

\sol{qm}{11} In units of $\mathrm{Ry} = 13.6\un{eV}$ and $u = a_0/b$, the Gaussian gives $\langle T\rangle = \tfrac32u^2$ and $\langle V\rangle = -(4/\sqrt\pi)u$. Minimizing $E(u)$ gives $u = 4/3\sqrt\pi$ and $E = -\tfrac{8}{3\pi}\,\mathrm{Ry} = -11.5\un{eV}$. The variational principle guarantees any trial energy lies above the true $-13.6\un{eV}$. The Gaussian is worse because it falls off too fast at large $r$ and is smooth at $r = 0$, where the $1/r$ potential demands a cusp; $e^{-r/a}$ has both features right, and is the exact answer.

\sol{qm}{12} Away from the origin, $\psi'' = \kappa^2\psi$ with $E = -\hbar^2\kappa^2/2m$, and normalizability with continuity gives $\psi = Ae^{-\kappa\abs x}$. Integrating the Schr\"odinger equation across $x = 0$: $-(\hbar^2/2m)[\psi'(0^+)-\psi'(0^-)] = \alpha\psi(0)$, i.e.\ $(\hbar^2/2m)(2\kappa A) = \alpha A$, so $\kappa = m\alpha/\hbar^2$ and $E = -m\alpha^2/2\hbar^2$. There is only one such $\kappa$, and odd functions vanish at the origin and never feel the delta, so there is exactly one bound state.

\sol{qm}{13} At $1\un{fm}$, $e^2/4\pi\epsilon_0r = 1.44\un{MeV\,fm}/1\un{fm} = 1.4\un{MeV}$. At $1.5\times10^7\un{K}$, $\kB T = 1.3\un{keV}$. With $E_G = 2(469\un{MeV})(\pi/137)^2 = 490\un{keV}$, the tunnelling factor is $e^{-\sqrt{490}} = 2\times10^{-10}$ at $1\un{keV}$ and $e^{-\sqrt{98}} = 5\times10^{-5}$ at $5\un{keV}$. The rate at energy $E$ goes as $e^{-E/\kB T}e^{-\sqrt{E_G/E}}$: the Boltzmann factor falls with $E$ while tunnelling rises, and their product peaks sharply at the \emph{Gamow peak}, $E_0 = (\kB T\sqrt{E_G}/2)^{2/3}\approx6\un{keV}$, about $4.5\,\kB T$. On top of the tiny tunnelling probability, two protons can only stick if one turns into a neutron by the weak interaction during their brief encounter, which almost never happens. A given proton in the Sun's core waits about ten billion years to fuse.

\sol{qm}{14} The first beam splitter puts amplitude $1/\sqrt2$ in each arm. With a bomb in arm~2: with probability $\tfrac12$ the photon is found there and the bomb explodes. Otherwise the photon is in arm~1, and since the other path is blocked there is no interference at the second beam splitter, which sends it to $A$ or $B$ with probability $\tfrac12$ each. So $25\%$ of the time $B$ clicks --- impossible without a bomb, since interference forbids it --- and the bomb is certified live without having absorbed a photon. ($25\%$ of the time $A$ clicks and the test is inconclusive; repeat it.) With many weak interrogation steps, the quantum Zeno effect keeps freezing the photon out of the bomb's path, pushing the success rate toward $100\%$.
